% !TEX root = ../main.tex
\chapter{ケーススタディ：DB スキーマ変更とクエリ保存}
\label{chap:db-schema-query-preservation}
\BookIndex{database schema}{DB スキーマ}
\BookIndex{くえりほぞん}{クエリ保存}
\BookIndex{join}{join}

\begin{chapterabstract}
本章では、DB スキーマ変更を「データの形を変える処理」としてではなく、「重要なクエリの意味を保つべき変換」として扱います。
テーブル分割、統合、カラム追加、正規化、非正規化は、実務では頻繁に起きます。
しかし、移行後のデータで既存のクエリが同じ答えを返すかどうかは、別途確認しなければなりません。
ここでは、実 DB エンジン全体を形式化するのではなく、小さな行単位モデルを Lean 4 で作り、旧注文行をユーザー部分と注文部分へ分けても表示用ビューが保存されることを証明します。
圏論の言葉では、Pullback を「共通キーによって整合するデータだけを合流させる考え方」として使います。
\end{chapterabstract}

\begin{learninggoals}
\item DB スキーマ変更で「保存すべきクエリ」を明示できます。
\item join を、共通キーに関する整合条件つきの合流として読めます。
\item Lean 4 で、旧テーブル、新テーブル、移行関数、クエリ関数を小さく定義できます。
\item 移行前後のクエリ結果一致を、等式として証明する形を理解できます。
\item Lean モデルで保証した範囲と、実 DB で別途確認すべき範囲を区別できます。
\end{learninggoals}

\section{問題設定}

DB マイグレーションでは、「スキーマ変更は終わった」「データも移した」といえても、それだけでは十分ではありません。
アプリケーションにとって重要なのは、多くの場合、テーブルの内部形ではなく、そこから得られる問い合わせ結果です。

たとえば、旧スキーマでは注文テーブルにユーザー名が重複して入っていたとします。
新スキーマでは、ユーザー情報を \texttt{users} テーブルに分け、注文情報を \texttt{orders} テーブルに置きます。
これは、正規化としては自然な変更です。
しかし、画面や API が必要としている「注文 ID、ユーザー名、金額」の一覧が、移行前後で変わっていないことは別の主張です。

\FigurePlaceholder{fig-ch32-schema-split.png}{0.92}{注文スキーマの分割}{fig:ch32-schema-split}

図\ref{fig:ch32-schema-split}では、左側に旧スキーマの横長の注文行があります。
右側には、ユーザー情報と注文情報へ分割された新スキーマがあります。
読者が見るべき点は、テーブルの形は変わっていても、画面や API に返したい「注文 ID、ユーザー名、金額」というビューは変わっていないという点です。

\section{ソフトウェア上の問題：スキーマ変更はクエリの意味を壊しうる}

「DB マイグレーションが成功した」という表現には、次の異なる条件が含まれえます。
\begin{itemize}
\item DDL を実行できました。
\item データ変換バッチが最後まで走りました。
\item 新アプリケーションが起動できました。
\item 主要クエリの結果が旧実装と一致しました。
\end{itemize}
本章で扱うのは、最後の意味です。
つまり、「アプリケーションが観測する答え」が保存されたかを扱います。

DB スキーマ変更には、少なくとも次のようなパターンがあります。

\begin{longtable}{p{0.25\linewidth}p{0.31\linewidth}p{0.34\linewidth}}
\toprule
変更 & 例 & 保存したい性質 \\
\midrule
テーブル分割 & \texttt{orders} から \texttt{users} を分ける & join 後の表示結果が同じ \\
テーブル統合 & \texttt{users} と \texttt{profiles} をまとめる & 旧クエリを projection で再現できる \\
カラム追加 & \texttt{status} を追加する & 旧クエリが追加カラムに依存しない \\
カラム削除 & 使われていない値を落とす & 削除値を参照するクエリがない \\
正規化 & 重複データをマスタへ移す & 外部キーで元の情報を復元できる \\
非正規化 & 集計済みカラムを追加する & 集計値が元データから計算した値と一致する \\
\bottomrule
\end{longtable}

この表の右列にある性質は、自然言語では「互換性がある」とまとめられがちです。
しかし、Lean で扱うには、より具体的に書く必要があります。
本章では、「旧 DB に対するクエリ」と「移行後 DB に対するクエリ」が同じ結果を返すという等式に落とします。

\section{クエリ結果が保存されるとは何か}

\begin{definitionbox}
本章でいうクエリ保存とは、旧 DB \(d\) に対して、次の形の等式が成り立つことです。
\[
  \mathit{queryNew}(\mathit{migrate}(d)) = \mathit{queryOld}(d)
\]
左辺は、旧 DB を新スキーマへ移行してから新クエリを実行した結果です。
右辺は、旧 DB に旧クエリを実行した結果です。
この二つが一致するとき、そのクエリに関しては移行が意味を保存しているといいます。
\end{definitionbox}

この定義は、すべてのクエリが保存されることを要求していません。
実務では、保存したいクエリを選ぶ必要があります。
たとえば、既存 API のレスポンスを支えるクエリ、帳票に使うクエリ、監査ログに出すクエリなどです。
保存すべきクエリを明示しないままスキーマ変更をレビューすると、「何が壊れていないなら安全なのか」が曖昧になります。

\begin{warningbox}
「スキーマが情報を失っていない」ことと、「特定のクエリ結果が保存される」ことは別です。
情報を失っていなくても、join 条件、重複行の扱い、NULL の扱い、並び順の扱いが変われば、特定のクエリ結果は変わりえます。
逆に、情報を一部落としていても、ある画面に必要なクエリだけは保存される場合もあります。
\end{warningbox}

この章では、証明対象を意図的に小さくします。
旧スキーマの 1 行を、新スキーマのユーザー行と注文行に分けます。
そして、旧行から直接作る表示用ビューと、新行を join して作る表示用ビューが一致することを示します。
その後、リスト全体に対して同じ性質を持ち上げます。

\FigurePlaceholder{fig-ch32-query-preservation-square.png}{0.88}{クエリ保存の可換図式}{fig:ch32-query-preservation-square}

図\ref{fig:ch32-query-preservation-square}は、本章の中心となる等式を図で表しています。
上から右へ行く道は、「旧 DB を移行してから新クエリを実行する」道です。
左から下へ行く道は、「旧 DB に旧クエリを実行する」道です。
二つの道の結果が同じなら、この図式は可換です。
ソフトウェアの言葉では、移行処理を挟んでも観測結果が変わりません。

\section{Pullback 的な整合制約}

DB join では、二つのテーブルから任意の行をペアにすればよいわけではありません。
ユーザー行と注文行を結合するなら、通常は \texttt{userId} が一致していなければなりません。
この「共通キーが一致する行だけを合流させる」という見方が、Pullback の直感に対応します。

\begin{definitionbox}
本章では、Pullback を完全な一般論として定義しません。
ここでは、Pullback 的な整合制約を次のように読みます。
\begin{quote}
二つのデータが共通の先へ写ったとき、同じ値になるものだけを組み合わせる。
\end{quote}
DB join であれば、ユーザー行と注文行が共通の \texttt{userId} を持つ場合だけ、それらを一つの表示結果へ合流させます。
\end{definitionbox}

SQL の inner join は条件を満たす行の組を作りますが、本章はその bag semantics を形式化しません。
\texttt{Consistent} という述語で、一組のユーザー部分と注文部分が同じ \texttt{userId} を指すことだけを表します。

\FigurePlaceholder{fig-ch32-pullback-join.png}{0.88}{共通キーによる整合}{fig:ch32-pullback-join}

図\ref{fig:ch32-pullback-join}では、ユーザー行から \texttt{userId} を取り出す矢印と、注文行から \texttt{userId} を取り出す矢印があります。
その二つが同じ値を返すときだけ、二つの行を組み合わせます。
この見方により、join は単なるペア作成ではなく、制約を満たす合流として読めます。

\section{Lean では行単位の分解モデルから始める}

Lean で扱うのは実 DB ではなく、旧注文一行から作ったユーザー部分と注文部分を同じ要素内に保持する行単位モデルです。
Users 表と Orders 表を独立に保持して join するモデルではありません。

まず、旧スキーマでは注文行にユーザー名が直接入っているとします。
新スキーマでは、ユーザー情報と注文情報を分けます。
表示用の \texttt{OrderView} は、アプリケーションが観測したいクエリ結果です。

\begin{listing}[htbp]
\begin{minted}{lean4}
structure OrderOld where
  orderId : Nat
  userId : Nat
  userName : String
  amount : Nat

deriving Repr, DecidableEq

structure UserRow where
  userId : Nat
  userName : String

deriving Repr, DecidableEq

structure OrderRow where
  orderId : Nat
  userId : Nat
  amount : Nat

deriving Repr, DecidableEq

structure SplitRow where
  user : UserRow
  order : OrderRow

deriving Repr, DecidableEq

structure OrderView where
  orderId : Nat
  userName : String
  amount : Nat

deriving Repr, DecidableEq
\end{minted}
\caption{\texttt{OrderOld}・\texttt{SplitRow}・\texttt{OrderView}}
\label{lst:ch32-schema-model}
\end{listing}

この段階では、まだ \texttt{SplitRow} が整合しているとは限りません。
たとえば、\texttt{user.userId} と \texttt{order.userId} が異なる値でも、型としては作れます。
そのため、後で \texttt{Consistent} という述語を導入します。

\section{一行の移行と一行のクエリ保存}

旧行を新スキーマへ分割する関数を定義します。
また、旧行から直接ビューを作る関数と、新しい分割行からビューを作る関数を定義します。
まず一行だけを見れば、証明したい性質は単純です。

\begin{listing}[htbp]
\begin{minted}{lean4}
def splitOrder (r : OrderOld) : SplitRow :=
  { user := { userId := r.userId, userName := r.userName },
    order := { orderId := r.orderId,
               userId := r.userId,
               amount := r.amount } }

def queryOldRow (r : OrderOld) : OrderView :=
  { orderId := r.orderId,
    userName := r.userName,
    amount := r.amount }

def querySplitRow (s : SplitRow) : OrderView :=
  { orderId := s.order.orderId,
    userName := s.user.userName,
    amount := s.order.amount }

theorem split_row_preserves_query (r : OrderOld) :
    querySplitRow (splitOrder r) = queryOldRow r := by
  cases r
  rfl
\end{minted}
\caption{\texttt{split\_row\_preserves\_query}}
\label{lst:ch32-row-preservation}
\end{listing}

証明は、\texttt{cases r} と \texttt{rfl} だけで終わります。
これは、移行関数とクエリ関数の定義を展開すると、両辺が同じレコードになるためです。
ここで重要なのは、証明が短いことそのものではありません。
重要なのは、「何を保存したいか」を \texttt{querySplitRow (splitOrder r) = queryOldRow r} という等式として明示した点です。

\begin{warningbox}
この一行の証明は、重複ユーザーの統合や外部キー制約の完全性を証明しているわけではありません。
証明しているのは、旧行から作ったユーザー部分と注文部分を使って表示用ビューを作ると、旧行から直接作ったビューと一致するという限定された性質です。
\end{warningbox}

\section{整合条件を型に持ち上げる}

次に、Pullback 的な整合制約を Lean の述語として表します。
新スキーマ側のユーザー行と注文行は、同じ \texttt{userId} を指している必要があります。
この条件を \texttt{Consistent} と呼びます。

\begin{listing}[htbp]
\begin{minted}{lean4}
def Consistent (s : SplitRow) : Prop :=
  s.user.userId = s.order.userId

def ConsistentSplitRow := { s : SplitRow // Consistent s }

theorem split_consistent (r : OrderOld) :
    Consistent (splitOrder r) := by
  cases r
  rfl

def splitOrderChecked (r : OrderOld) : ConsistentSplitRow :=
  Subtype.mk (splitOrder r) (split_consistent r)

def queryNewRow (s : ConsistentSplitRow) : OrderView :=
  { orderId := s.val.order.orderId,
    userName := s.val.user.userName,
    amount := s.val.order.amount }
\end{minted}
\caption{\texttt{ConsistentSplitRow}}
\label{lst:ch32-consistency}
\end{listing}

\texttt{ConsistentSplitRow} は、値 \texttt{s : SplitRow} と、その値が \texttt{Consistent s} を満たす証明をまとめた型です。
これは、\ChapterRef{chap:ch10_invariants}で扱った「条件を満たす値」を使う考え方の応用です。
DB の言葉への類比では、単なる二つの行の組ではなく、key が一致している証拠を一緒に持つ行です。
ただし、この部分型そのものは実 DB の外部キー制約ではありません。

この型を使うと、新スキーマ側のクエリは「整合している行」にだけ適用されます。
実 DB では、制約や join 条件がこの役割を担います。
Lean の小さなモデルでは、述語と部分型がその役割を担います。

\section{リスト全体の行単位モデルに持ち上げる}

実務の DB は、一行ではありません。
そこで、旧 DB を旧行のリスト、検証用の移行後モデルを整合済み分割行のリストとして表します。
移行は \texttt{List.map splitOrderChecked} です。
旧クエリと新クエリも、それぞれ一行クエリをリスト全体へ写す関数として定義します。

\begin{listing}[htbp]
\begin{minted}{lean4}
def DBV1 := List OrderOld

def DBV2 := List ConsistentSplitRow

def migrateDB (db : DBV1) : DBV2 :=
  List.map splitOrderChecked db

def queryOld (db : DBV1) : List OrderView :=
  List.map queryOldRow db

def queryNew (db : DBV2) : List OrderView :=
  List.map queryNewRow db

theorem checked_row_preserves_query (r : OrderOld) :
    queryNewRow (splitOrderChecked r) = queryOldRow r := by
  cases r
  rfl

theorem query_preserved (db : DBV1) :
    queryNew (migrateDB db) = queryOld db := by
  induction db with
  | nil =>
      rfl
  | cons x xs ih =>
      simp [migrateDB, queryNew, queryOld,
            checked_row_preserves_query]
\end{minted}
\caption{\texttt{query\_preserved}}
\label{lst:ch32-list-preservation}
\end{listing}

この定理が、本章の中心です。
ただし、証明対象は行単位モデルの \texttt{List.map} であり、独立した Users 表と Orders 表に対する SQL join ではありません。
\[
  \texttt{queryNew (migrateDB db) = queryOld db}
\]
という形で、任意の旧 DB \texttt{db} について、移行後に新クエリを実行した結果と、移行前に旧クエリを実行した結果が一致することを述べています。

証明では、\texttt{induction db} を使っています。
空リストの場合は自明です。
先頭行 \texttt{x} と残り \texttt{xs} の場合は、一行の保存性 \texttt{checked\_row\_preserves\_query} と、帰納法の仮定 \texttt{ih} を使って整理します。
これは、\ChapterRef{chap:ch30-list-migration-functor-laws}で扱った「一件の移行をリスト全体へ持ち上げる」考え方と同じです。

\section{サンプルを実行して結果を見る}

証明は、すべての入力に対する主張です。
それとは別に、\texttt{\#eval} で小さな例を実行すると、モデルの意味を確認しやすくなります。

\begin{listing}[htbp]
\begin{minted}{lean4}
def sampleOld : DBV1 :=
  [ { orderId := 1001, userId := 10,
      userName := "Ada", amount := 120 },
    { orderId := 1002, userId := 20,
      userName := "Grace", amount := 250 } ]

#eval queryOld sampleOld  -- => [{ orderId := 1001, userName := "Ada", amount := 120 }, { orderId := 1002, userName := "Grace", amount := 250 }]
#eval queryNew (migrateDB sampleOld)  -- => [{ orderId := 1001, userName := "Ada", amount := 120 }, { orderId := 1002, userName := "Grace", amount := 250 }]
\end{minted}
\caption{\texttt{sampleOld}}
\label{lst:ch32-sample-eval}
\end{listing}

この二つの \texttt{\#eval} は、同じ表示用ビューのリストを返します。
ただし、\texttt{\#eval} は具体例の確認であり、証明ではありません。
本章の本体は、\texttt{query\_preserved} という定理です。
\texttt{\#eval} は、読者が定義を読む補助として使います。

\section{この定理が保証していること}

\texttt{query\_preserved} が保証していることは、次のように限定して読むべきです。

\begin{itemize}
\item 旧 DB を \texttt{migrateDB} で新 DB へ移した場合、\texttt{queryNew} の結果は \texttt{queryOld} の結果と一致します。
\item 新 DB の各行は、\texttt{ConsistentSplitRow} によって \texttt{userId} の整合性を持ちます。
\item 任意の長さの旧 DB リストについて、この性質が成り立ちます。
\end{itemize}

この定理は、実務レビューで次のような説明に置き換えられます。

\begin{quote}
この検証モデルでは、旧注文行からユーザー情報と注文情報を分割しても、注文一覧画面が使う \texttt{orderId, userName, amount} の結果は変わらないことを証明している。
\end{quote}

このように説明できると、レビューの焦点が明確になります。
「移行は安全です」という広すぎる主張ではなく、「このクエリについて、このモデルの前提のもとで結果が保存されます」といえます。

\section{DB 検証モデルを実装へ適用する}

この定理を実 DB へ適用するには、Lean モデルの外側にある条件も検査対象として対応付けます。
設計レビューでは、少なくとも次の項目を通常のテスト、DB 制約、移行リハーサルで確認します。

\begin{warningbox}
Lean モデルと併用する検査項目は次のとおりです。
\begin{itemize}
\item SQL の実行計画、インデックス、ロック、トランザクション分離レベル
\item NULL、重複行、照合順序、タイムゾーン、文字コードを含む DB エンジン固有の挙動
\item 本番データが移行スクリプトの前提を満たすことと、途中失敗時の復旧手順
\item 同一 key の重複、孤児注文、SQL の bag semantics における重複数と結果順序
\end{itemize}
\end{warningbox}

Users 表に同じ \texttt{userId} の行が二つあれば、inner join は一件の注文から複数結果を作りえます。
対応するユーザーがなければ、結果は 0 件になります。
現在の \texttt{DBV2} は各旧行にユーザー部分を埋め込むため、これらの境界事例を表現しません。
実 DB へ適用する際は、キー一意性、外部キー、NULL 方針、set/bag semantics を別の前提やテストとして固定します。

Lean モデルはクエリ保存の核を検査し、DB 制約とテストは実エンジン上の前提を検査します。
両者の対応を設計文書に残すことで、証明結果を実装レビューへ適用できます。

図\ref{fig:ch32-lean-boundary}では、中央の行変換、整合述語、クエリ等式と、周辺の SQL 方言、DB 制約、移行運用、性能検証を分けます。

\FigurePlaceholder{fig-ch32-lean-boundary.png}{0.90}{DB 検証モデルの境界}{fig:ch32-lean-boundary}

\section{実務へ持ち込むときの手順}

この章の考え方を実務へ持ち込むなら、次の順で進めるとよいでしょう。

\begin{enumerate}
\item 保存したいクエリを一つ選びます。
\item そのクエリの出力型を、表示用ビューや API レスポンス型として書きます。
\item 旧スキーマからその出力型を作る関数を定義します。
\item 新スキーマから同じ出力型を作る関数を定義します。
\item 移行関数を定義します。
\item \texttt{queryNew (migrate old) = queryOld old} の形で定理を書きます。
\item 必要なら、外部キー一致や well-formed 条件を述語または部分型として入れます。
\item 証明した範囲と、実 DB で別途確認する範囲を設計文書に書きます。
\end{enumerate}

重要なのは、最初からすべてのクエリを対象にしないことです。
まず一つの重要クエリを選び、移行仕様の核を Lean に移します。
その小さな成功例が、次のクエリや次のマイグレーションへ広げる足場になります。

\section{レビューで使うチェックリスト}

最後に、DB スキーマ変更をレビューするときの確認項目をまとめます。
これは、Lean コードを書く前にも使えます。

保存したいクエリが、自然言語ではなく出力型として明示されているかを確認します。

旧スキーマから出力型を作る関数があるかを確認します。

新スキーマから同じ出力型を作る関数があるかを確認します。

移行関数が、旧スキーマから新スキーマへの変換として定義されているかを確認します。

証明したい性質が \texttt{queryNew (migrate db) = queryOld db} の形になっているかを確認します。

join に必要な整合条件、たとえば外部キー一致が述語として明示されているかを確認します。

証明モデルが扱わない SQL 固有の挙動が、別途テストやレビューで扱われているかを確認します。

\section{本章のまとめ}

本章では、DB スキーマ変更をクエリ保存の問題として扱いました。
旧スキーマと新スキーマを小さな Lean の構造体として表し、移行関数とクエリ関数を定義しました。
そして、旧 DB に旧クエリを実行した結果と、移行後 DB に新クエリを実行した結果が一致することを証明しました。

スキーマ変更の安全性は、「どのクエリ結果を保存するか」を明示して初めて検査しやすくなります。

Pullback 的な見方は、join を共通キーに関する整合条件つきの合流として理解する助けになります。

Lean では、実 DB 全体ではなく、小さな有限モデルで移行仕様の核を表します。

\texttt{queryNew (migrateDB db) = queryOld db} は、クエリ保存を表す基本形です。

Lean の証明は強力ですが、SQL 方言、NULL、トランザクション、性能、本番データ品質は別途扱う必要があります。
